\section{System Design and Architecture}
\label{sec:design}

\subsection{Overview}
\proj{} is structured as a \emph{pipeline} that converts an incoming PDF request
into searchable DOM text. A PDF navigation is intercepted and redirected to a
custom viewer; the viewer renders each page, checks whether it paints raster images or lacks
extractable text, recovers image text with OCR where needed, normalizes coordinates,
and injects an invisible text layer that the browser's find feature can search.
Figure~\ref{fig:pipeline} shows the end-to-end flow.

\begin{figure}[p]
  \centering
  \vspace*{6mm}
  \begin{tikzpicture}[
      %% Generous vertical gap between every node
      node distance=14mm,
      %% Main-column boxes
      box/.style={rectangle, rounded corners=4pt, draw=accent!70, fill=accent!6,
                  text width=6.4cm, align=center, minimum height=13mm,
                  font=\small, inner sep=7pt},
      %% OCR-branch boxes (slightly shaded to mark the conditional path)
      side/.style={rectangle, rounded corners=4pt, draw=accent!80, fill=accent!12,
                   text width=4.6cm, align=center, minimum height=13mm,
                   font=\small, inner sep=7pt},
      dec/.style={diamond, aspect=3.0, draw=accent!70, fill=accent!6,
                  align=center, font=\small, inner sep=3pt},
      io/.style={rectangle, draw=rulegray, fill=codebg, text width=6.4cm,
                 align=center, font=\small, minimum height=10mm, inner sep=5pt},
      arr/.style={-{Stealth[length=2.6mm]}, draw=accent!80, thick},
      lbl/.style={font=\scriptsize\itshape, fill=white, inner sep=2pt},
    ]

    %% ── Left (main) column ──────────────────────────────────────────────────
    \node[io]  (req)    {PDF request
                         (\code{http} / \code{https} / \code{file})};

    \node[box, below=of req]    (sw)
      {Service worker redirects to custom viewer\\
       via \code{declarativeNetRequest}};

    \node[box, below=of sw]     (render)
      {PDF.js renders page to \code{<canvas>}\\
       and extracts native text items};

    \node[dec, below=14mm of render, text width=3.2cm] (dec)
      {Images present\\or no text?};

    \node[box, below=46mm of dec] (norm)
      {Normalize word coordinates;\\
       cache results in IndexedDB};

    \node[box, below=of norm]   (inject)
      {Inject invisible \code{<span>} text layer\\
       positioned over the canvas};

    \node[io,  below=of inject] (search)
      {Searchable via custom find bar
       ($\mathrm{Ctrl}{+}\mathrm{F}$)};

    %% ── Right (OCR) branch ──────────────────────────────────────────────────
    \node[side, right=26mm of dec] (ocr)
      {Tesseract.js OCR\\
       (Web Worker pool)};

    \node[side, below=14mm of ocr] (merge)
      {Merge \& de-duplicate\\
       native + OCR words};

    %% ── Arrows ──────────────────────────────────────────────────────────────
    \draw[arr] (req)    -- (sw);
    \draw[arr] (sw)     -- (render);
    \draw[arr] (render) -- (dec);

    %% "no" path: straight down the left column, skips OCR
    \draw[arr] (dec.south)
      -- node[lbl, left=4pt] {no --- text-only page}
      (norm.north);

    %% "yes" path: right to OCR branch
    \draw[arr] (dec.east)
      -- node[lbl, above=3pt] {yes --- images found}
      (ocr.west);

    \draw[arr] (ocr) -- (merge);

    %% merge rejoins main column via a clean L-shaped connector
    \draw[arr] (merge.south) -- ++(0,-10mm) -| (norm.east);

    \draw[arr] (norm)   -- (inject);
    \draw[arr] (inject) -- (search);
  \end{tikzpicture}
  \vspace*{4mm}
  \caption{End-to-end processing pipeline. Each page is rendered and its native text
           extracted. If the page contains raster image operators \emph{or} has no
           extractable text (right branch), Tesseract.js runs OCR on a high-resolution
           offscreen canvas and the results are merged with any native text, dropping
           duplicates. Text-only pages skip OCR and pass straight to normalization
           (left path). Both paths converge at coordinate normalization before the
           invisible text layer is injected over the canvas.}
  \label{fig:pipeline}
\end{figure}

\subsection{Components}
\begin{description}[leftmargin=1.4em, style=nextline]
  \item[Request interception] A service worker registers a dynamic
        \code{declarativeNetRequest} rule on install that matches URLs ending in
        \code{.pdf} and redirects them to \code{viewer.html}, preserving the
        original URL as a query parameter.
  \item[Custom viewer] A self-contained HTML/JS page that owns rendering, OCR
        orchestration, the text layer, and the find UI.
  \item[Rendering] PDF.js rasterizes each page to a canvas, lazily and at the
        device pixel ratio for crisp output.
  \item[Text detection] For each page the viewer asks PDF.js for its text content
        and operator list to decide whether OCR is required.
  \item[OCR] A pool of Tesseract.js Web Workers recognizes text on image pages,
        returning words with bounding boxes.
  \item[Normalization \& caching] Word coordinates are normalized to the unit
        square and cached in IndexedDB so results are reusable across zoom levels
        and sessions.
  \item[Text-layer injection] Invisible, individually-scaled \code{<span>}
        elements are positioned over the canvas.
  \item[Search] A custom find bar searches the injected text and draws highlight
        overlays.
\end{description}

\subsection{The Coordinate-Mapping Problem}
The hardest part of the system is geometric. Text originates in up to four
distinct coordinate spaces, and the injected overlay must align with the rendered
page in all of them simultaneously:

\begin{center}
\begin{tikzpicture}[
    sp/.style={rectangle, rounded corners=2pt, draw=accent!70, fill=accent!6,
               align=center, font=\small\ttfamily, minimum height=9mm,
               inner sep=4pt},
    a/.style={-{Stealth[length=2.2mm]}, draw=accent!80, thick}]
  \node[sp] (pdf) {PDF space\\\footnotesize\rmfamily points, $\uparrow$ origin};
  \node[sp, right=7mm of pdf] (vp) {Viewport\\\footnotesize\rmfamily PDF.js scale};
  \node[sp, right=7mm of vp] (cv) {Canvas\\\footnotesize\rmfamily device px};
  \node[sp, right=7mm of cv] (dom) {DOM overlay\\\footnotesize\rmfamily CSS px, $\downarrow$ origin};
  \draw[a] (pdf) -- (vp);
  \draw[a] (vp) -- (cv);
  \draw[a] (cv) -- (dom);
\end{tikzpicture}
\end{center}

PDF space uses typographic points with a bottom-left origin; the PDF.js viewport
applies a scale (and any page rotation); the canvas backing store is further
scaled by the device pixel ratio for high-DPI sharpness; and the DOM overlay uses
CSS pixels with a top-left origin. A glyph recovered or extracted in one space
must be drawn at the corresponding location in the overlay, or search highlights
and text selection drift away from the visible characters.

For \emph{native} PDF text, PDF.js supplies, per item, an affine transform
$M_{\text{item}}$ and the viewport's transform $M_{\text{vp}}$ (which already
folds in page rotation). The glyph's full matrix in device space is their product
\begin{equation}
  M \;=\; M_{\text{vp}}\, M_{\text{item}},
  \qquad
  M=\begin{bmatrix} a & c & e\\ b & d & f \\ 0 & 0 & 1\end{bmatrix},
  \label{eq:matmul}
\end{equation}
from which the baseline angle and device-space font height follow directly:
\begin{equation}
  \theta = \operatorname{atan2}(b,\,a),
  \qquad
  h = \sqrt{c^{2}+d^{2}} .
  \label{eq:angle}
\end{equation}
Because rotation falls out of the matrix in Eq.~\eqref{eq:angle}, rotated pages
and rotated text are handled without special cases. The span's top-left corner is
obtained by stepping from the baseline origin $(e,f)$ up by the ascent
$h_a \approx 0.8\,h$ along the text direction.

For \emph{OCR} text, Tesseract reports pixel bounding boxes in the
high-resolution OCR canvas. These are normalized to the unit square by dividing
by the OCR canvas dimensions,
\begin{equation}
  \tilde{x} = \frac{x}{W_{\text{ocr}}},\qquad
  \tilde{y} = \frac{y}{H_{\text{ocr}}},
  \label{eq:normalize}
\end{equation}
so the cached result is independent of the resolution at which OCR happened. At
draw time they are denormalized into the text layer's CSS-pixel space using the
\emph{viewport} dimensions---not the DPR-scaled canvas dimensions---so the
invisible text and the visible page share one coordinate system.

\subsection{Key Design Decisions}
Four decisions shape the rest of the system and are justified where they appear
in Section~\ref{sec:impl}:

\begin{enumerate}
  \item \textbf{Never modify the original PDF.} The system is a strictly
        read-only render plus an overlay: the source file's bytes and format are
        left untouched, and all recovered text lives in a separate, invisible DOM
        layer. This constraint, fixed from the outset, rules out any
        rewrite-the-PDF approach and is the reason searchability is added
        \emph{around} the document rather than \emph{into} it.
  \item \textbf{Inject real text instead of intercepting search.} This reuses the
        browser's indexing, matching, and (optionally) highlighting, and keeps
        recovered text selectable and copyable.
  \item \textbf{Separate the OCR image from the display image.} Display fidelity
        and OCR legibility have different requirements, so the page is rendered
        twice: once for the screen at the device pixel ratio, and once offscreen
        at high resolution purely for recognition.
  \item \textbf{Lazy rendering with background indexing.} Pages render as they
        approach the viewport for responsiveness, while a separate idle-time pass
        guarantees that off-screen pages still become searchable.
\end{enumerate}
